%% ------------------------------------------------------------------
%% clip_intro_paragraph.tex
%% One paragraph for the Introduction, after the paragraph that formalises
%% CLIP through the rule (g(z),f(u)) >= theta and before "In this paper we
%% present ...".  Bib entries: clip_bib_entries.bib; radford2021learning
%% and udandarao2024no already exist.
%% ------------------------------------------------------------------

That this rule is a good classifier can be checked directly, and the check tells us what has to be studied. We used the open-source implementation of CLIP \cite{radford2021learning} provided by OpenCLIP \cite{ilharco2021openclip}, with the ViT-B/32 image encoder (a Vision Transformer \cite{dosovitskiy2021vit} with $32\times32$ patches) and the checkpoint \texttt{laion2b\_s34b\_b79k} trained on the LAION-2B dataset \cite{schuhmann2022laion,cherti2023reproducible}, to embed the $10{,}000$ test images of CIFAR-100 \cite{krizhevsky2009cifar} into $\Real^{512}$ and, with the prompt ``a photo of a \{\}'', the names of its $100$ classes; the code and data are available at \url{https://github.com/Mirkes/OneShotSeparabilityTheorem}. For every class $z$, the halfspace $\{x:(g(z),x)\ge\theta\}$ separates the images of class $z$ from the images of all other classes with ROC AUC $0.98$ on average and $0.91$ at worst over the $100$ classes (Section~\ref{sec:clip-examples}). Thus the text embedding $g(z)$ is a good normal vector for a halfspace that captures the class, and it is constructed from one object only---one description of the class---without any image of it. This is the situation studied in this paper: a hyperplane built from a single representative of a class, which must leave the representatives of all other classes on the other side. CLIP learns an embedding in which, for the captions it has seen, images rank correctly against sampled alternatives, and then one hopes that the halfspaces of $m$ new labels will not capture each other's representatives. That hope is a geometric property of the learned embedding rather than a term of the loss, and it is the one-shot separability of the set of representatives defined in Section~\ref{sec:one-shot:general}. The capacity of zero-shot classification---how many classes can be told apart by such a rule---is therefore governed by the one-shot separability of the embeddings of class descriptions, and in Section~\ref{sec:clip-examples} we compare its behaviour on these data with the theory below: the rate at which the probability that one embedding violates the halfspace of another decays with the dimension, and the number of simultaneously usable labels, which follows the resulting bound until it is limited by the data rather than by the dimension.
